% 18 November 2026. Prepared for Edward Baker with OpenAI Codex LLM
% assistance. Actual inherited serving model and effort unavailable.
\section{One initialized divisor on a common time interval}
\label{sec:common-time}

The compatible polynomial construction in Section~\ref{sec:finite}
admits an explicit approximation theorem on a fixed time interval.
The essential point is that the polynomial is initialized once: it is
not replaced by a new truncation of the canonical divisor at each time.
We give the proof from \cite{BakerCM47}, assessed within its fixed-input
scope in \cite{BakerCMReview40}. Its quantitative conclusion is an
$O(\log N/N)$ field error with specified constants. A useful numerical
realization additionally needs those constants and finite input errors
to satisfy the contour tests in Section~\ref{sec:counted-transfer}.

\subsection{Initialization and the omitted opposite-pair tail}

Fix $I=[t_0,t_1]\subset[0,1/2]$, put $\Delta=t_1-t_0$, and retain
the complete canonical source and initial count assumptions of
Section~\ref{sec:framework}. Let $h_0=H(t_0,0)>0$ and
$H_0(z)=H(t_0,z)$. Choose one representative $\rho_\nu$ for each
opposite pair of zero \emph{occurrences} at $t_0$. Thus multiplicity
$m$ at each of $\rho,-\rho$ contributes $m$ factors. The multiset of
squares $\rho_\nu^2$ is conjugation invariant, and
\begin{equation}
 H_0(z)=h_0\prod_\nu(1-z^2/\rho_\nu^2).
 \label{cta:initial-product}
\end{equation}
For a divisor cutoff $N\ge1$, retain every modulus tie, opposite pair,
conjugate and multiplicity, and define
\begin{equation}
 \begin{split}
 J_N&=\#\{\nu:|\rho_\nu|\le N\},\\
 P_N(t_0,z)&=h_0\prod_{|\rho_\nu|\le N}(1-z^2/\rho_\nu^2),\\
 P_N(t_0+\tau,z)&=\exp(-\tau\partial_z^2)P_N(t_0,z),
       \qquad 0\le\tau\le\Delta,\\
 E_N&=\sum_{|\rho_\nu|>N}|\rho_\nu|^{-2}
     =\frac12\sum_{\substack{\rho\text{ an occurrence}\\|\rho|>N}}
                    |\rho|^{-2}.
 \end{split}
 \label{cta:initialization}
\end{equation}
The exponential terminates. Throughout these two sections $N$ is a
divisor radius; the moment-head order $P$ in Section~\ref{sec:finite}
is a different parameter. No global strip or reality assertion for
the evolved polynomial is part of the construction.

\begin{lemma}[A common absolute-product envelope]
\label{cta:envelope}
Set $K=5\Cu$ and supply the finite near-origin data
\begin{equation}
 q_{<1}=\#\{\nu:|\rho_\nu|<1\},\qquad
 C_{<1}=\prod_{|\rho_\nu|<1}\max(1,|\rho_\nu|^{-2}).
 \label{cta:near-data}
\end{equation}
With an empty product equal to one, define
\begin{equation}
 g_0=h_0C_{<1},\qquad A=2q_{<1}+8K,\qquad
 G(r)=g_0\exp\{A(1+r)\log(2+r)\}.
 \label{cta:envelope-definition}
\end{equation}
Then the full absolute product and every partial absolute product obey
\begin{equation}
 h_0\prod_\nu(1+r^2/|\rho_\nu|^2)\le G(r),\qquad r\ge0,
 \label{cta:absolute-product}
\end{equation}
and, for $N\ge1$,
\begin{equation}
 J_N\le KN\log(2+N),\qquad
 E_N\le\frac{2K\{\log(2+N)+1\}}N.
 \label{cta:tail-bound}
\end{equation}
The envelope is independent of $N$.
\end{lemma}

\begin{proof}
Let $A_*(x)=\#\{\nu:|\rho_\nu|\le x\}$. Covering the projected
interval $[-x,x]$ by integer unit bands and applying
\eqref{eq:count-safe} gives, for $x\ge1$,
\[
 A_*(x)\le\tfrac12(2x+3)\Cu\log(3+x)
          \le Kx\log(2+x).
\]
The factor $1/2$ counts opposite occurrences exactly. It uses neither a
lower population bound nor a zero-count asymptotic. The factors of
modulus below one contribute at most $C_{<1}(1+r)^{2q_{<1}}$.
For all remaining factors, the positive integral identity for
$\log(1+r^2/s^2)$ and interchange of nonnegative sums give
\begin{equation}
 \sum_{|\rho_\nu|\ge1}\log(1+r^2/|\rho_\nu|^2)
 \le2Kr^2\int_1^\infty\frac{\log(2+u)}{u^2+r^2}\dd u.
 \label{cta:layer-cake}
\end{equation}
Factors of modulus one are included by this integral identity.
For $r\ge1$, split the integral at $r$. The lower and upper pieces
are at most $2Kr\log(2+r)$ and $2Kr\{\log(2+r)+1\}$, since
\begin{equation}
 \int_b^\infty\frac{\log(2+u)}{u^2}\dd u
 =\frac{\log(2+b)}b+\frac12\log(1+2/b)
 \le\frac{\log(2+b)+1}b.
 \label{cta:tail-integral}
\end{equation}
Their sum is at most $6Kr\log(2+r)$. For $0\le r\le1$, the right
side of \eqref{cta:layer-cake} is at most
$2Kr^2(\log3+1)<5Kr^2$. These bounds and the near-origin factors
prove \eqref{cta:absolute-product} with the larger common constant
in \eqref{cta:envelope-definition}.

Finally, again interchanging nonnegative terms,
\[
 E_N=\int_N^\infty2u^{-3}
              \#\{\nu:N<|\rho_\nu|\le u\}\dd u
 \le2K\int_N^\infty\frac{\log(2+u)}{u^2}\dd u.
\]
Equation~\eqref{cta:tail-integral} proves the tail assertion; the
count assertion is the bound for $A_*(N)$.
\end{proof}

\begin{lemma}[Initial field error without division by a factor]
\label{cta:telescoping}
For every $z\in\C$ and $N\ge1$,
\begin{equation}
 |H_0(z)-P_N(t_0,z)|\le |z|^2G(|z|)E_N.
 \label{cta:initial-error}
\end{equation}
The alternative estimate
$|H_0-P_N(t_0,\cdot)|\le G(|z|)(\exp(|z|^2E_N)-1)$ also holds.
\end{lemma}

\begin{proof}
At a larger complete finite cutoff, add the omitted factors one at a
time. The change on adding $1-z^2/\rho_\nu^2$ is $-z^2/\rho_\nu^2$
times a partial product. Its absolute value, including $h_0$, is at
most $|z|^2|\rho_\nu|^{-2}G(|z|)$ by
Lemma~\ref{cta:envelope}. Summing and passing to the convergent
canonical product proves \eqref{cta:initial-error}. No possibly
vanishing factor was divided out, so the estimate holds also at simple
and multiple zeros. For the alternative bound, factor the retained
product from the omitted product and use
$|\prod(1+a_\nu)-1|\le\exp(\sum|a_\nu|)-1$.
\end{proof}

\subsection{An evaluated Gaussian propagation constant}

\begin{theorem}[Common-time field and spatial-jet approximation]
\label{cta:field-theorem}
Keep the fixed initial inputs above. Define
\begin{equation}
 \begin{gathered}
 \lambda=\frac1{8(1+\Delta)},\qquad
 \varepsilon_\lambda=\frac{\lambda}{2A},\qquad
 b_\lambda=2\varepsilon_\lambda-1-\log\varepsilon_\lambda,\\
 B_\lambda=A b_\lambda+
       \frac{A^2(\varepsilon_\lambda+b_\lambda)^2}{2\lambda},
 \qquad M_\lambda=g_0\e^{B_\lambda},\\
 \mathcal C(R,\Delta)=M_\lambda\e^{2\lambda R^2}
       \left\{2R^2\sqrt{1+\Delta}
                     +4\Delta(1+\Delta)^{3/2}\right\}.
 \end{gathered}
 \label{cta:gaussian-constant}
\end{equation}
For every fixed $R\ge0$ and every $N\ge1$,
\begin{equation}
 \sup_{\substack{t\in I\\|z|\le R}}
       |H(t,z)-P_N(t,z)|\le\mathcal C(R,\Delta)E_N.
 \label{cta:field-error}
\end{equation}
For each padding $p>0$ and $0\le j\le3$, a common jet bound is
\begin{equation}
 \sup_{\substack{t\in I\\|z|\le R}}
       |\partial_z^j(H-P_N)(t,z)|\le C_jE_N,
 \qquad C_j=\frac{j!}{p^j}\mathcal C(R+p,\Delta).
 \label{cta:jet-error}
\end{equation}
For $j=0$ the constant in \eqref{cta:field-error} is slightly sharper.
All these errors are $O(\log N/N)$ at the fixed inputs. No small-tail
condition is needed for the field estimate.
\end{theorem}

\begin{proof}
For a standard normal real variable $X$, the terminating Taylor series
and $\mathbb E X^{2k}=(2k)!/(2^kk!)$ give
\begin{equation}
 P_N(t_0+\tau,z)
   =\mathbb E P_N(t_0,z+\ii\sqrt{2\tau}X).
 \label{cta:polynomial-gaussian}
\end{equation}
The term of order $2k$ is precisely
$(-\tau)^k\partial_z^{2k}P_N(t_0,z)/k!$, with the required negative
heat sign. For the canonical field this formula must be justified
separately from the literal source. For $u\ge0$,
\[
 \mathbb E|\cos((z+\ii\sqrt{2\tau}X)u)|
 \le \e^{|\Im z|u}\mathbb E\e^{\sqrt{2\tau}|X|u}
 \le2\e^{|\Im z|u+\tau u^2}.
\]
After multiplication by $\e^{t_0u^2}\Phi(u)$, the source decay in
\eqref{eq:source} gives an integrable majorant, uniformly on the fixed
compact spatial set and $0\le\tau\le\Delta$. Fubini and the Gaussian
exponential expectation therefore prove
\begin{equation}
 H(t_0+\tau,z)=\mathbb E H_0(z+\ii\sqrt{2\tau}X).
 \label{cta:canonical-gaussian}
\end{equation}
This does not assert backward-heat existence for an arbitrary entire
initial product.

To evaluate a dominating constant, use
$\log y\le\varepsilon_\lambda y-1-\log\varepsilon_\lambda$.
It implies $\log(2+r)\le\varepsilon_\lambda r+b_\lambda$ and hence
\[
 A(1+r)\log(2+r)
 \le\tfrac\lambda2r^2+A(\varepsilon_\lambda+b_\lambda)r+A b_\lambda
 \le\lambda r^2+B_\lambda.
\]
The last step is $ur\le\lambda r^2/2+u^2/(2\lambda)$.
Thus $G(r)\le M_\lambda\e^{\lambda r^2}$.
For $|z|\le R$ and $0\le\tau\le\Delta$, set
$r=R+\sqrt{2\Delta}|X|$. Then
$r^2\le2R^2+4\Delta X^2$ and
$1-8\lambda\Delta=(1+\Delta)^{-1}$. Direct Gaussian integration gives
\[
 \mathbb E\e^{4\lambda\Delta X^2}=\sqrt{1+\Delta},\qquad
 \mathbb E X^2\e^{4\lambda\Delta X^2}=(1+\Delta)^{3/2}.
\]
Consequently $\mathbb E[r^2G(r)]\le\mathcal C(R,\Delta)$.
Apply Lemma~\ref{cta:telescoping} inside
\eqref{cta:polynomial-gaussian}--\eqref{cta:canonical-gaussian} to obtain
\eqref{cta:field-error}. This evaluates the whole Gaussian path,
including its tails. Finally, Cauchy's formula on a circle of radius
$p$ about each $|z|\le R$ proves \eqref{cta:jet-error}.
\end{proof}

Both fields satisfy the same heat equation, so the estimates for orders
two and three also pay $\partial_t(H-P_N)$ and
$\partial_t\partial_z(H-P_N)$. The theorem concerns fixed height,
time interval and derivative order. Its constants can be large; their
explicit finiteness proves neither affordable acquisition nor a uniform
height or simultaneous order/cutoff limit.

\subsection{The finite coefficient realization}

Write $J=J_N$ and $P_N(t,z)=\sum_{k=0}^J a_{2k}(t)z^{2k}$.
Coefficient comparison in the heat equation gives
\begin{equation}
 a_{2k}'=-(2k+2)(2k+1)a_{2k+2},\qquad a_{2J}'=0,
 \label{cta:coefficient-flow}
\end{equation}
and repeated differentiation, or the terminating exponential, gives
\begin{equation}
 a_{2k}(t_0+\tau)=\sum_{r=0}^{J-k}
 \frac{(-\tau)^r}{r!}\frac{(2k+2r)!}{(2k)!}a_{2k+2r}(t_0).
 \label{cta:coefficient-solution}
\end{equation}
This is a closed system of $J+1$ real coefficients. Real even symmetry,
degree and the leading coefficient are preserved, including at mergers.
All higher moments of this polynomial are determined by its divisor;
they have not been assigned independent zero values above a moment head.

The required finite input is the complete initial divisor through $N$,
its ties and multiplicities, and its normalization, or equivalent
coefficient data. The complete count/tail input, near-origin factors and
source identity are additional analytical inputs. Lemma~\ref{cta:envelope}
bounds the number of required coefficients; it does not acquire them.
Section~\ref{sec:counted-transfer} pays their uncertainty and the
independent contour information needed for the intended observables.
